\documentclass[12pt,a4paper]{article}
\usepackage[utf8]{inputenc}
\usepackage[spanish]{babel}
\usepackage{geometry}
\usepackage{enumitem}
\usepackage{titlesec}
\usepackage{xcolor}

\geometry{a4paper, margin=1in}

\titleformat{\section}{\normalfont\Large\bfseries\color{blue}}{}{0pt}{}
\titleformat{\subsection}{\normalfont\large\bfseries\color{darkgray}}{}{0pt}{}

\title{Distribución de Tareas - Semana 1 \\ \large Curso: Seguridad Informática}
\author{Prof. César Rodríguez}
\date{\today}

\begin{document}

\maketitle

\section{Introducción}
Para el inicio del proyecto, los 13 estudiantes se dividirán en 4 grupos de trabajo. Cada estudiante es responsable de implementar una funcionalidad específica en el archivo de su grupo dentro del paquete \texttt{modulos/}.

\section{Grupo 1: Reconocimiento DNS (3 Estudiantes)}
\textbf{Archivo:} \texttt{modulos/dns\_recon.py}

\begin{itemize}[label=$\bullet$]
    \item \textbf{Estudiante 1:} Implementar \texttt{get\_a\_records}. Consultar registros \textbf{A} y \textbf{AAAA}.
    \item \textbf{Estudiante 2:} Implementar \texttt{get\_mx\_ns\_records}. Consultar registros \textbf{MX} y \textbf{NS}.
    \item \textbf{Estudiante 3:} Implementar \texttt{get\_txt\_soa\_records}. Consultar registros \textbf{TXT} y \textbf{SOA}.
\end{itemize}

\section{Grupo 2: OSINT y Huella Digital (3 Estudiantes)}
\textbf{Archivo:} \texttt{modulos/osint.py}

\begin{itemize}[label=$\bullet$]
    \item \textbf{Estudiante 1:} Implementar consulta \textbf{WHOIS} (usando \texttt{python-whois}).
    \item \textbf{Estudiante 2:} Implementar búsqueda de \textbf{subdominios} mediante Google Dorks.
    \item \textbf{Estudiante 3:} Implementar detección de \textbf{archivos expuestos} (\texttt{.log}, \texttt{.env}, \texttt{.sql}).
\end{itemize}

\section{Grupo 3: Descubrimiento de Hosts (3 Estudiantes)}
\textbf{Archivo:} \texttt{modulos/discovery.py}

\begin{itemize}[label=$\bullet$]
    \item \textbf{Estudiante 1:} Implementar \textbf{Ping Sweep} básico usando paquetes ICMP.
    \item \textbf{Estudiante 2:} Crear lógica para procesamiento de \textbf{rangos de red} (CIDR a lista de IPs).
    \item \textbf{Estudiante 3:} Implementar detección mediante \textbf{TCP SYN Ping} a puertos comunes.
\end{itemize}

\section{Grupo 4: Escaneo de Puertos y Servicios (4 Estudiantes)}
\textbf{Archivo:} \texttt{modulos/scanning.py}

\begin{itemize}[label=$\bullet$]
    \item \textbf{Estudiante 1:} Implementar escáner \textbf{TCP Connect} básico.
    \item \textbf{Estudiante 2:} Implementar lógica de \textbf{concurrencia} (uso de \texttt{threading} o \texttt{concurrent.futures}).
    \item \textbf{Estudiante 3:} Implementar escáner \textbf{UDP} y manejo de respuestas ICMP.
    \item \textbf{Estudiante 4:} Implementar \textbf{Banner Grabbing} inicial para identificación de servicios.
\end{itemize}

\section{Instrucciones de Entrega}
\begin{enumerate}
    \item \textbf{Contrato de Datos:} Toda función debe retornar un diccionario que cumpla con \texttt{docs/schema\_resultados.json}.
    \item \textbf{No Bloqueo:} El programa principal no debe detenerse por errores en los módulos. Use \texttt{try-except}.
    \item \textbf{Git:} Trabaje únicamente en el archivo de su grupo. No modifique \texttt{auditoria.py}.
    \item \textbf{Documentación:} Use Docstrings estilo Google y Type Hinting.
\end{enumerate}

\section{Anexo: Guía de Estándares de Código}

\subsection{Type Hinting (Tipado de Datos)}
Permite declarar qué tipo de datos espera una función y qué devolverá, ayudando al editor a detectar errores.
\begin{itemize}
    \item \textbf{Argumentos:} \texttt{nombre\_variable: tipo} (ej: \texttt{domain: str}).
    \item \textbf{Retorno:} \texttt{-> tipo} (ej: \texttt{-> dict}).
\end{itemize}

\subsection{Docstrings (Estilo Google)}
Es el formato estándar para documentar la lógica de la función. Debe incluir:
\begin{itemize}
    \item \textbf{Args:} Lista de parámetros, su tipo y una breve descripción.
    \item \textbf{Returns:} Descripción del valor de retorno y su estructura.
\end{itemize}

\textbf{Ejemplo de referencia para los estudiantes:}
\begin{verbatim}
def mi_funcion(dominio: str) -> dict:
    """
    Realiza una consulta específica al dominio.
    Args:
        dominio (str): El nombre de dominio a investigar.
    Returns:
        dict: Resultados siguiendo el esquema JSON definido.
    """
    pass
\end{verbatim}

\vfill
\centering
\textit{Recuerde consultar el manual 'Hackers 6' disponible en el repositorio.}

\end{document}